\documentclass[11pt,a4paper,headsepline]{scrartcl}
\usepackage[utf8]{inputenc}
\usepackage[german]{babel}
\usepackage{amsmath}
\usepackage{amsfonts}
\usepackage{amssymb}
\usepackage{tikz}
\usepackage[left=2cm,right=2cm,top=2cm,bottom=2cm]{geometry}
\usepackage[utf8]{inputenc}
\usepackage{graphicx}
\usepackage{pgfplots}
\usepackage{fancyhdr}
\usepackage[scaled]{helvet}
\usepackage{paralist}

\usetikzlibrary{calc}
\newtheorem{aufgabe}{Aufgabe}


\renewcommand*\familydefault{\sfdefault}
\renewcommand{\arraystretch}{1.1}

\title{\"Ubung L\"angsdynamik}
\date{}
\makeatletter
\let\Title\@title
\let\Author\@author
\makeatother
\pagestyle{fancy}
\fancyhead[L]{Prof. Dr. Raphael Pfaff}
\fancyhead[R]{84111: Schienenfahrzeugtechnik}
\fancyfoot[L]{Datei: \jobname}
\fancyfoot[R]{Datum: \today
\begin{picture}(0,0)(0,0)\put(.5,0){\includegraphics[height=5cm]{fh_logo}}\end{picture}}

\begin{document}
\maketitle
\thispagestyle{fancy}
\pagestyle{fancy}
\vspace{-2cm}

\begin{aufgabe}[Railway Challenge]
\label{Ta:Acceleration}
Eine Parkbahn-Lokomotive soll gem\"a{\ss} Lastenheft ausgelegt werden. Bestimmen Sie:
\begin{itemize}
	\item Die f\"ur die Acceleration Challenge optimale Masse innerhalb der Grenzen der Spezifikation f\"ur einen konstanten Reibwert $\mu = 0{,}15$.
	\item Den Energiebedarf f\"ur den Dauerbetrieb gem\"a{\ss} Spezifikation.
	\begin{itemize}
		\item Rechnen Sie in eine Masse des Energiespeichers um f\"ur folgende Energiedichten: \vspace{.2cm} \\ 
		\centering
		\begin{tabular}{|c|c|}
		\hline
		Wasserstoff & $1{,}19 \frac{\mathrm{MJ}}{\mathrm{kg}}$ \\ \hline
		Bleiakkumulator & $0{,}11 \frac{\mathrm{MJ}}{\mathrm{kg}}$ \\ \hline
		Mi-Mh-Akku & $0{,}28 \frac{\mathrm{MJ}}{\mathrm{kg}}$ \\ \hline
		Methan & $50 \frac{\mathrm{MJ}}{\mathrm{kg}}$ \\ \hline
		Diesel & $43 \frac{\mathrm{MJ}}{\mathrm{kg}}$ \\ \hline
		\end{tabular}
		\end{itemize}
		\item Bestimmen Sie die Energie der vorgeschriebenen rekuperativen Bremsung und die daraus zu erreichende Geschwindigkeit abh\"angig vom Wirkungsgrad. Weiterhin bestimmen Sie die zu erreichende Fahrstrecke unter der Annahme von einem Fahrwiderstand von $1\%$ der Gewichtskraft abh\"angig vom Wirkungsgrad.
\end{itemize}
\end{aufgabe}
\newpage
\begin{aufgabe}
Eine Lokomotive der BR 143 zieht einen Wagenzug. Die technischen Daten der Fahrzeuge sind:
\begin{itemize}
	\item Triebfahrzeug:
	\begin{itemize}
		\item Masse $m_{L} = 82\, \mathrm{ t}$
		\item Rotierende Masse $m_{DL} = 16\, \mathrm{ t}$
	\end{itemize}
	\item Wagenzug:
	\begin{itemize}
		\item Masse $m_{W} = 500\, \mathrm{t}$
		\item Rotierende Masse $m_{DW} = 24\, \mathrm{t}$
	\end{itemize}
\end{itemize}

\pgfplotsset{width=10 cm}
\begin{center}
	\begin{tikzpicture}
		\begin{axis}[ylabel={$F_{T}/kN$},
		xlabel = {$v/ms^{-1}$}
		%legend, %style={at={(1.5,1)}},
		%anchor=north, legend columns=3},
		%xtick=data
		]]
		\addplot coordinates
			{(0,244) (2.78, 223) (5.55,207) (8.33, 196) (11.11, 189) (13.89, 183) (16.67, 178) (19.44, 174) (22.22, 170) (25, 168) (26.11, 167) (27.78, 144) (29.17, 129) (30.56, 117) (31.94, 107) (33.33, 100)};
			\addplot coordinates
			{(0,11.2) (2.78, 11.5) (5.55, 12.5) (8.33, 13.3) (11.11, 14.8) (13.89, 16.7) (16.67, 18.9) (19.44, 21.6) (22.22, 24.6) (25, 28) (26.11,29.4 ) (27.78, 31.7) (29.17, 33.8) (30.56, 35.9) (31.94, 38.1) (33.33, 40.5)};
			%\addplot coordinates %2%-Gefaelle
			%{(0,111.2) (2.78, 111.5) (5.55, 112.5) (8.33, 113.3) (11.11, 114.8) (13.89, 116.7) (16.67, 118.9) (19.44, 121.6) (22.22, 124.6) (25, 128) (26.11,129.4 ) (27.78, 131.7) (29.17, 133.8) (30.56, 135.9) (31.94, 138.1) (33.33, 140.5)};
		\legend{$F_{T}(v)$, $F_{WZ}(v)$};
		\end{axis}
\end{tikzpicture}
\end{center}

	\begin{itemize}[a)]
		\item Zeichnen Sie die Widerstandskurven des Zugverbands (bestehend aus Lokomotive und Wagenzug) f\"ur Streckenneigungen $i_{k} = (2, 4) \%$, $k = 1,2$ in das $F-v$-Diagramm ein. Der Fahrwiderstand des Triebfahrzeugs ist zu vernachl\"assigen.
			\item Bestimmen Sie die H\"ochstgeschwindigkeiten $v_{max, k}$ in den jeweiligen Streckenneigungen.
		\item Bestimmen Sie das Beschleunigungsverm\"ogen des Zugverbands in der Ebene und in 1\% Streckenneigung f\"ur $v = 90\, \mathrm{km/h}$.
		\item Bestimmen Sie die Ausnutzung des Rad-Schiene-Kraftschluss f\"ur das Erreichen des maximalen Beschleunigungsverm\"ogens.
		\end{itemize}
\end{aufgabe}
\end{document}